% MG01. Inserción: después de exc:aw y su interpretación elíptica;
% antes de exc:codigo. Requiere L_1 y la carta A_W ya construidos.
\subsubsection{Neutralidad dual y orientación de la primera frontera conjunta}
\label{mg01:frontera}

La frontera de prolongación reúne las tres biografías regionales mediante
una relación que conserva el eje, el agregado lateral, la saturación y
la ocupación de las columnas. En la coordenatización de vectores fila ya fijada,
los últimos bloques de longitud seis de las palabras de longitud treinta
son
\[
 u_4^\pi=110221,\qquad u_4^e=102011,\qquad
 u_4^\varphi=010200.
\]
Los superíndices conservan las tres regiones de
\eqref{eq:palabras-regionales}; cada palabra se interpreta como un
elemento de \(\mathbb F_3^6\), con representantes enteros en
\(\{0,1,2\}\). La relación de frontera transporta el eje mediante
\begin{equation}
 z:=u_5^\varphi=u_4^e=102011,
 \qquad a:=u_4^\varphi A_WL_1^3=111101.
 \label{mg01:eje-contraste}
\end{equation}
Aquí \(a\) es el contraste orientado entre los dos canales laterales y
el eje. El agregado lateral correspondiente es
\begin{equation}
 x+y=210112,\qquad
 q:=x+y+z=012120,\qquad a=x+y-z=111101
 \quad\text{en }\mathbb F_3^6,
 \label{mg01:agregados}
\end{equation}
donde \(x=u_5^\pi\) e \(y=u_5^e\). La suma y el contraste tienen
lecturas duales distintas: \(qA_W=012000\) y \(aA_W=120210\).

Sea \(B\in M_{3,6}(\mathbb F_3)\) la matriz de filas \(x,y,z\).
Su levantamiento canónico permite definir, columna a columna,
\[
 c_j=\left\lfloor\frac{x_j+y_j+z_j}{3}\right\rfloor,
 \qquad
 o_j=\mathbf1_{x_j\ne0}+\mathbf1_{y_j\ne0}+\mathbf1_{z_j\ne0}.
\]
% BEGIN R28_BOUNDARY_OCCUPANCY_INPUT
% BEGIN R28_BOUNDARY_OCCUPANCY
% Insertar en MG01 tras la definición de c_j,o_j, antes de su valor.
La ley de la frontera de profundidad seis fija \(c_j=1\). Con el eje \(z\) y el
agregado \(q\) ya determinados, las sumas enteras laterales son
\[
 (x_j+y_j)_{j=1}^6=(2,4,3,4,4,2).
\]
La regla de ocupación mínima selecciona, entre sus descomposiciones con
\(x_j,y_j\in\{0,1,2\}\), las columnas con menor número de entradas
no nulas, contando también \(z_j\). Para suma dos, las parejas extremas
\((0,2),(2,0)\) ocupan una posición lateral; para suma tres, las parejas
\((1,2),(2,1)\) ocupan dos; para suma cuatro, sólo existe \((2,2)\).
Al añadir la ocupación del eje se obtiene \(o=(2,2,3,2,3,2)\).
En esta frontera, con \(s=[x+y]_3\), la misma regla se expresa como
\[
 o_j=2+\mathbf1_{\{z_j\ne0,\ s_j\ne2\}}.
\]
% END R28_BOUNDARY_OCCUPANCY

% END R28_BOUNDARY_OCCUPANCY_INPUT

En esta frontera, la saturación y la ocupación son
\begin{equation}
 c=111111,\qquad o=223232.
 \label{mg01:saturacion}
\end{equation}
Por tanto, la igualdad entera de cada columna es
\(x_j+y_j+z_j=q_j+3\). El acarreo \(c\) registra el levantamiento
entero de esta frontera; su dominio es el de las seis columnas presentes.

\begin{proposition}[Ocho distribuciones de una frontera saturada]
\label{mg01:ocho}
Las condiciones \eqref{mg01:eje-contraste}--\eqref{mg01:saturacion}
dejan exactamente ocho matrices. Sus alternativas columnares son
\begin{equation}
\begin{array}{c|c|c|c}
 j&z_j&x_j+y_j\ \text{en }\mathbb Z&(x_j,y_j)\\\hline
 1&1&2&(0,2),(2,0)\\
 2&0&4&(2,2)\\
 3&2&3&(1,2),(2,1)\\
 4&0&4&(2,2)\\
 5&1&4&(2,2)\\
 6&1&2&(0,2),(2,0)
\end{array}
\label{mg01:tabla-columnas}
\end{equation}
\end{proposition}
\begin{proof}
La tercera fila queda fijada por el transporte del eje. La igualdad
entera \(x_j+y_j=q_j+3-z_j\), junto con \(o_j\), da las alternativas
de la tabla. Las columnas primera, tercera y sexta admiten dos elecciones
independientes; las otras tres admiten una. El producto de sus cardinales
es \(2\cdot1\cdot2\cdot1\cdot1\cdot2=8\), y cada elección conserva
todas las igualdades declaradas.
\end{proof}

La neutralidad se evalúa en la coordenatización dual. Con
\(\mathbf1=(1,1,1,1,1,1)^{\mathsf T}\), la matriz de
\eqref{exc:aw} determina
\begin{equation}
 \omega=A_W\mathbf1=(2,1,1,1,1,1)^{\mathsf T},
 \qquad A_W\omega=-\mathbf1.
 \label{mg01:unidad-dual}
\end{equation}
La condición de cierre dual de la frontera es
\begin{equation}
 (BA_W)\mathbf1=B\omega=0\quad\text{en }\mathbb F_3^3.
 \label{mg01:neutralidad}
\end{equation}
Esta condición actúa sobre la colección de fronteras saturadas. La relación
cuadrática de \(A_W\) determina el transporte de la unidad dual; la
neutralidad especifica qué fronteras satisfacen el cierre correspondiente.

\begin{proposition}[Pareja dual y frontera orientada]
\label{mg01:dos-orientacion}
La neutralidad \eqref{mg01:neutralidad} deja exactamente
\begin{equation}
 B_-=
 \begin{pmatrix}0&2&1&2&2&2\\2&2&2&2&2&0\\1&0&2&0&1&1\end{pmatrix},
 \qquad
 B_+=
 \begin{pmatrix}2&2&2&2&2&0\\0&2&1&2&2&2\\1&0&2&0&1&1\end{pmatrix}.
 \label{mg01:pareja}
\end{equation}
La involución que intercambia clausura y propagación permuta \(B_-\)
y \(B_+\). En la coordenatización semiabierta con la orientación APP fijada para
los canales, la frontera canónica es \(R_{36}=B_+\).
\end{proposition}
\begin{proof}
Escribimos las tres elecciones binarias de la primera fila como
\[
 x=(2\varepsilon_1,2,1+\varepsilon_3,2,2,2\varepsilon_6),
 \qquad \varepsilon_1,\varepsilon_3,\varepsilon_6\in\{0,1\}.
\]
La condición \(x\omega=0\) equivale a
\[
 1+\varepsilon_1+\varepsilon_3-\varepsilon_6\equiv0\pmod3.
\]
Las posibilidades son \((0,0,1)\) y \((1,1,0)\). La fila del eje
satisface \(z\omega=0\); el agregado lateral también tiene evaluación
dual nula. Por consiguiente, en ambos casos \(y\omega=0\). Resultan
exactamente las dos matrices de \eqref{mg01:pareja}.

La elección final aplica la orientación APP de la coordenatización semiabierta a
esta pareja especular. En dicha coordenatización, las representaciones ordenadas son
\[
 B_-\longmapsto(789,048,894),\qquad
 B_+\longmapsto(793,045,894).
\]
El representante de la orientación fijada es \(B_+\). Esta operación
de selección conserva el orden de las tres regiones y la convención de
frontera de sus cilindros. La reducción algebraica de ocho matrices a
dos y la elección de representante orientado constituyen, así, dos
etapas de la misma construcción.
\end{proof}

Las cargas visibles distinguen esas etapas con precisión:
\begin{equation}
 B_-\mathbf1=(0,1,2)^{\mathsf T},\qquad
 B_+\mathbf1=(1,0,2)^{\mathsf T},\qquad
 B_\pm\omega=0.
 \label{mg01:cargas}
\end{equation}
La carga visible y la neutralidad dual son evaluaciones sobre vectores
distintos. La frontera seleccionada conserva tres filas de seis trits,
con índices ambientales \((726,215,301)\), y proporciona el estado
de empalme \(R_{36}\) de la prolongación conjunta. Su memoria incluye
el eje transportado, la saturación, la incidencia de columnas y la
orientación que distingue las hojas de la pareja.
